Visual classifiers trained with standard augmentation pipelines tend to exploit
spurious correlations: background statistics, texture patterns, or scene-level
cues that happen to co-occur with class labels but do not reflect object
structure.  Conventional augmentation strategies, including random crops, colour
jitter, and policy-based methods such as RandAugment and AutoAugment, are
constructed without reference to what the model has actually learned, and so
offer no mechanism for targeting the specific shortcuts a given classifier has
acquired.

We present Intelligent Coarse Dropout (\textbf{ICD}) and its complement
Anti-ICD (\textbf{AICD}), two data augmentation methods that condition spatial
masking on the model's own class activation maps.  Standard coarse dropout
(cutout) hides rectangular regions chosen at random, erasing background as
readily as object.  ICD and AICD make this saliency-aware: they split the image
into a grid of tiles, score each tile by its mean saliency, and selectively
mask tiles by a percentile threshold.  ICD masks the highest-saliency
tiles, the regions the model already relies on, so that the classifier must
recruit additional cues.  AICD masks the lowest-saliency tiles, perturbing the
background while leaving the discriminative region intact.  Masked tiles are
not filled with a hard black box but with one of several soft strategies
(blurred original, local mean, noise, or a constant), which preserve different
aspects of local context.

We describe the formal construction of both masks, the family of fill
strategies, the role of the tile size and threshold hyperparameters, and a
reference implementation in the open-source BNNR library.  This paper
introduces the methods; a quantitative evaluation against standard augmentation
baselines is left to a companion study.
